% !TeX root = ../main.tex
% Versão de revisão: corpo das entradas preservado.
% \nota{...} são acréscimos propostos, já visíveis no PDF.
% NOTA SUGERIDA marca uma intervenção para você escrever; não aparece no PDF.
% CONFERIR indica dado ou localização ainda a verificar na edição consultada.
% As três chamadas \notaarquivo originais continuam apontando para seus arquivos.
% Para usar a lista ao final deste arquivo, substitua o conteúdo de
% referencias.tex por \referenciasdodiario. main.tex conserva seu título
% de seção e continua chamando % Lista MANUAL em ordem alfabética, sem numeração, conforme MLA 9/RDQ.
% Inclua apenas as obras realmente citadas e confira os dados da edição usada.
% Títulos significativos em maiúsculas iniciais; livros em itálico;
% artigos/capítulos entre aspas; não informar cidade para livros.
% Referências abreviadas ficam NAS NOTAS, não em chamadas autor-data.

\referencia{Agência Câmara de Notícias. \enquote{Joice Hasselmann Diz que Deixa a Liderança no Congresso com Sentimento de Dever Cumprido}. \textit{Câmara dos Deputados}, 17 out. 2019. Disponível em: \url{https://www.camara.leg.br/noticias/601207-joice-hasselmann-diz-que-deixa-a-lideranca-no-congresso-com-sentimento-de-dever-cumprido/}. Acesso em: 21 nov. 2026.}

\referencia{Barros, Douglas Rodrigues. \enquote{Contra Nego Bispo}. \textit{Blog da Boitempo}, 21 jan. 2026. Disponível em: \url{https://blogdaboitempo.com.br/2026/01/21/contra-nego-bispo/}. Acesso em: 20 nov. 2026.}

\referencia{Carneiro da Cunha, Manuela. \enquote{\enquote{Cultura} e cultura: conhecimentos tradicionais e direitos intelectuais}. \textit{Cultura com Aspas e Outros Ensaios}. Ubu Editora, 2017.}

\referencia{Foucault, Michel. \enquote{As heterotopias}. \textit{O corpo utópico, as heterotopias}, n-1 edições, 2013.}

\referencia{Foucault, Michel. \textit{As Palavras e as coisas}, Martins Fontes, 1999.}

\referencia{Haraway, Donna. \enquote{Pensamento tentacular: Antropoceno, Capitaloceno, Chthuluceno}. \textit{Ficar com o problema: Fazer parentes no Chthluceno}, n-1 edições, 2023. p 59-116.}

\referencia{Krenak, Ailton. \textit{Futuro Ancestral}. Companhia das Letras, 2022.}

\referencia{Latour, Bruno. \enquote{How to Talk About the Body? The Normative Dimension of Science Studies}. \textit{Body \& Society}, vol.~10, n.~2--3, 2004, pp.~205--229. \url{https://doi.org/10.1177/1357034X04042943}.}

\referencia{Latour, Bruno. \textit{Jamais Fomos Modernos: Ensaio de Antropologia Simétrica}. Traduzido por Carlos Irineu da Costa, Editora 34, 2019.}

\referencia{Mauss, Marcel; Beuchat, Henri. \enquote{Ensaio sobre as variações sazonais das sociedades esquimós}. \textit{Sociologia e Antropologia}, escrito por Marcel Mauss, Cosac Naify, 2015. p. 423-501.}

\referencia{R7. \enquote{Ex-Líder do Governo Bolsonaro, Joice Hasselmann Declara Voto em Lula no 2º Turno}. \textit{R7 Notícias}, 23 out. 2026. Disponível em: \url{https://noticias.r7.com/eleicoes/2026/ex-lider-do-governo-bolsonaro-joice-hasselmann-declara-voto-em-lula-no-2-turno-23102026/}. Acesso em: 21 nov. 2026.}

\referencia{Santos, Antônio Bispo dos. \enquote{Dois e Dois São Dois: Renato Noguera e Nêgo Bispo}. Entrevista concedida a Renato Noguera. \textit{Amarello}, n.~40, 2022. Disponível em: \url{https://amarello.com.br/2022/04/cultura/dois-e-dois-sao-dois-renato-noguera-e-nego-bispo/}. Acesso em: 20 nov. 2026.}

\referencia{Santos, Antônio Bispo dos. \enquote{Entrevista: Nego Bispo}. Entrevista concedida a Gabriel Galo em 17 e 18 de junho de 2020. \textit{Papo de Galo}, n.~3, 2020. Disponível em: \url{https://papodegalo.com.br/entrevista-nego-bispo/}. Acesso em: 20 nov. 2026.}

\referencia{Santos, Antônio Bispo dos. \textit{A Terra Dá, a Terra Quer}. Piseagrama e Ubu Editora, 2023.}

\referencia{Santos, Gabriel. \enquote{Uma Crítica Marxista ao Pensamento Decolonial e a Nego Bispo}. \textit{Esquerda Online}, 10 fev. 2026. Disponível em: \url{https://esquerdaonline.com.br/2026/02/10/uma-critica-marxista-ao-pensamento-decolonial-e-a-nego-bispo/}. Acesso em: 20 nov. 2026.}. A lista não se duplica.

\entrada{16 de outubro de 2026}{Pensar a comunicação}

Esta pesquisa começa no dia 2 de abril de 2024 -- data que marca o dia internacional de conscientização do autismo -- quando vi  numa rede social uma publicação condenando os memes com características de transtornos psíquicos compartilhados por quem vive-com eles.
% NOTA SUGERIDA [cena de pesquisa]: Identificar em termos gerais o que o meme fazia circular e o que a postagem recusava. Basta uma cena ou exemplo não identificável; não é preciso reproduzir a postagem.

Sendo eu uma pessoa autista, me incomodei consideravelmente com aquela provocação naquele dia específico. Memificar nossas características em comum é um jeito de falarmos de nossos corpos por vias que operam desvios frente ao utilitarismo do vocabulário diagnóstico; é um jeito de constituir-com nossos corpos identidades mediadas positivamente, criativamente, e não pelo intermédio do defeito. Na lembrança daquela ocasião, enxergo-com os olhos de hoje os erros que Latour atribui aos pós-modernos: tentam desconstruir os diagnosticados apagando junto às doenças as diferenças mediadas hoje pelos diagnósticos; a grita despatologizante de nada ajuda aos patologizados podando suas diferenças.\nota{Latour, \textit{Jamais Fomos Modernos}. O paralelo com a despatologização é uma leitura deste diário.}

Decidi que escreveria um texto respondendo à postagem, buscando reintegrar as diferenças ao discurso sobre o corpo, e comecei a pesquisa que o articularia. Logo a proposta evoluiu para um trabalho maior que almejei então transformar num projeto de iniciação científica\notaarquivo{notas/ic.tex}: eu queria falar de autismo nas veredas entre as ciências médicas --- patologizantes --- e as vertentes das humanidades críticas à biologização da vida. Eu queria manter as diferenças sem enquadrá-las como defeitos, resolver a controvérsia adicionando-lhe uma dimensão: os quase-sujeitos em questão.

Latour descreve a Constituição dos modernos como um princípio de proliferação de quase-objetos -- mediadores que são deixados pelo caminho da estabilização de sujeitos e objetos.\nota{Latour, \textit{Jamais Fomos Modernos}, cap. 2.} Digo que somos enquadrados como quase-sujeitos na controvérsia em questão porque, de um ao outro lado, traduzimos circunstancialmente uma polarização ética na qual nossas diferenças específicas são tolhidas de qualquer valor, quando não da própria voz. Aqueles que patologizam as diferenças, ao menos têm a decência de reconhecê-las; enquanto os que negam rotular os corpos com defeitos, assim rotulam o mundo, apagando as diferenças de habituação entre corpos distintos.

A questão foi ficando ainda mais interessante pois incorporaria também uma dimensão política. A diferenciação pelo enquadramento da doença ou da deficiência medeia o acesso a certos direitos, além de articular discursos que mobilizam solidariedades. Como fazer, portanto, para trocar o figurino de defeito de nossas diferenças por outro -- de valência positiva -- sem acarretar desatenção aos modos específicos como elas nos proporcionam existir?
% NOTA SUGERIDA [escopo]: Se a passagem sobre direitos se tornar argumento jurídico, delimitar país e direito e citar a norma correspondente; por ora, é a tensão política que importa.

Busquei a orientação do professor Pedro Peixoto por ter sido o único em meu instituto em cujo currículo encontrei alguma menção ao autismo. Seu trabalho no campo dos estudos sociais das ciências e técnicas o coloca em companhia frequente com o pensamento de Bruno Latour -- intelectual que eu até então desconhecia. Pensando-com Latour e Peixoto compreendi que faltava uma dimensão ao enquadramento moderno da comunicação: para que as diferenças patologizadas sejam positivadas, o significado de uma mensagem deve ser o resultado de uma performance, não um atributo rígido da própria mensagem\notaarquivo{notas/peixoto.tex}.
% NOTA SUGERIDA [anonimato]: Para a cópia de submissão, retirar o nome do orientador e as pistas identificadoras em notas/ic.tex e notas/peixoto.tex. A segunda nota aparece incompleta no PDF.

O significado é um acontecimento no emissor, e outro no receptor --- que constrói um significado à parte, articulável em alguma medida com o primeiro. Tal é a medida da confluência comunicacional. Com isso minha pesquisa tomou outro rumo: o de compor enquadramentos para a comunicação e a ciência (enquanto processo genérico pelo qual se veria a ter ciência de algo) de modo a agregar o ruído como parte da mensagem, e não como ruído.

\entrada{17 de outubro de 2026}{Pensar o Antropoceno}

O dossiê para o qual escrevo pede pensar o Antropoceno, e minha primeira entrada neste diário de escrita -- escrita que é o próprio diário -- não parece ter se aproximado nem introduzido o tema. Penso em dois caminhos que realizam este contato: um que  segue a cronologia linear que descrevi anteriormente; outro que rompe com ela e espirala a partir da inadequação dos enquadramentos modernos à comunicação e à ciência.

Seguirei agora pelo segundo. Proponho pensar o Antropoceno a partir dessa inadequação. Encontro mais agência diante das emergências naturais e sociais do presente nomeando um \enquote{Antropoceno} --- com aspas --- que revela no processo  de sua nomenclatura a elipse dos quase-objetos quase-sujeitos, e na presença das aspas um estrangeirismo do termo. Demarcar uma época geológica pelos  impactos humanos sobre o sistema planetário figura seres não-humanos -- como montanhas, rios, oceanos, florestas, a atmosfera, o clima, entre tantos outros -- como objetos, como se de fato tivessem sido meros recursos e os ditos impactos humanos resultassem da agência humana, e não de uma associação de agências humanas e não-humanas; já a demarcação de um nome com as aspas ao mesmo tempo que nos livra de falar como se fôssemos antropocênicos, reintegra a dimensão que estes retiram da mensagem; a visão simultânea de ambas as demarcações proporciona agir-comunicar-com quase-objetos quase-sujeitos  mais articuladamente.\nota{Em 2000, Crutzen e Stoermer associaram a hipótese de um início do Antropoceno no fim do século XVIII ao aumento de gases de efeito estufa registrado no gelo e à máquina a vapor (\enquote{The Anthropocene}, p. 18). A proposta formal do Anthropocene Working Group escolheu o nível correspondente a 1952 em Crawford Lake, onde o aumento de plutônio fornece um marcador para correlacionar estratos na escala da Grande Aceleração. O voto contra a formalização, ratificado em 2024, deixou em aberto o uso mais amplo do termo (Subcommission on Quaternary Stratigraphy, \enquote{Working Group}; International Union of Geological Sciences e International Commission on Stratigraphy, \enquote{Joint Statement}).\par A mudança do marco me interessa porque torna visível o trabalho necessário para dizer que lugares distantes atravessaram a mesma fronteira temporal. Da máquina a vapor aos testes nucleares, mudam a escala das associações e os meios materiais que permitem reconhecê-las. As aspas começam, para mim, onde pergunto quais dessas associações o nome torna comunicáveis e quais responsabilidades ele pode diluir.}

Foquemos agora na primeira demarcação: a nomenclatura do \enquote{Antropoceno}. Em sua escolha pelos antropocênicos revela-se a autovisão destes, assim como seu caráter precipitado. Enquadram-se como uma força geológica, mas o fazem a partir da observação de efeitos geológicos não previstos nem desejados da aplicação desta força; enquadram-se como uma força, mas desgovernada, imprevisível; o fazem pois não conseguem enxergar agenciamento e delegação entre os seres não-humanos, porque se precipitam e tomam como objetos seres que funcionam como sujeitos, pois são incapazes de enquadrar não-humanos enquanto tais.
% NOTA SUGERIDA [escala e agência]: Articular o truque de Arquimedes de Latour à escala dos coletivos da máquina a vapor e dos testes nucleares. Distinguir agência distribuída de responsabilidade histórica desigual e nomear o coletivo europeu-modernizante-ocidental sem convertê-lo numa humanidade indiferenciada.

Se a atmosfera for figurada como sujeito, demarcar  transformações climáticas  nocivas  aos humanos como  antropogências seria como atribuir aos senhores de escravos mortos na Revolução Haitiana a autoria maior deste evento, e não àqueles humanos que reagem à  escravidão alçados por  seus próprios interesses.  Ao retirar-lhes a dignidade, os humanos convidam os não-humanos às animosidades. Nos quilombos, conforme ensina Nego Bispo, compartilha-se a dignidade com o outro e com os não-humanos, convidando todos às confluências.\nota{Santos, \textit{A Terra Dá, a Terra Quer}.}
% NOTA SUGERIDA [limites da analogia]: Explicitar o que se compara (a invisibilização da agência) e o que não se equipara (a ação histórica de pessoas escravizadas e as dinâmicas atmosféricas). Se mantida como mais do que analogia, documentar a passagem; considerar outro exemplo.

Mesmo que seja para denunciar os efeitos  adversos das agências humanas para a habitabilidade planetária aos humanos, a nomenclatura do \enquote{Antropoceno} se fundamenta no mesmo monoteísmo que orienta tais agências.

\entrada{18 de outubro de 2026}{Retornar ao que foi comunicado}

Ontem parei a escrita logo após uma menção a Nego Bispo. Nem sequer terminei de escrever o último parágrafo(em algum momento irei falar da  estrutura  discursiva que orienta minha escrita ainda nesta publicação). Comecei a escrevê-lo  mas  interrompi a escrita pois o mestre quilombola chamava a uma releitura de \enquote{A terra dá, a terra quer} antes de retomar a composição do texto com uma presença maior de Bispo.\nota{Santos, \textit{A Terra Dá, a Terra Quer}.}

Leitura inspiradora e potente, cuja leveza espanta. Concluí o reencontro em  menos de 24 horas. Decidi em meio ao acontecimento presentear o livro a duas pessoas às quais suas palavras me levaram todo o tempo: uma pela confluência com a ancestralidade, outra pela confluência com a vida no mato.

Me intrigou notar como Bispo constitui riquíssimas encruzilhadas. Dos modos de vida rurais com os saberes ancestrais; mas também entre os modos de se relacionar com a terra e os modos de se relacionar com os humanos (que, para ele, consistem nos povos cosmofóbicos e monoteístas das cidades).  Esta segunda encruza me levou a pensar em dois textos: um do antropólogo francês Marcel Mauss, outro da portuguesa abrasileirada Manuela Carneiro da Cunha\notaarquivo{notas/manuela_brasileira.tex}.

Em seu \enquote{Ensaio sobre as variações sazonais das sociedades esquimós} Mauss  me deixou fascinado com os povos Innuit (nomenclatura contracolonial daqueles que o francês  chamou de esquimós --  \enquote{comedores de foca} em sua língua nativa). O  estudo que elabora traduz dos modos de vida  Innuit -- na língua nativa, \enquote{o povo} -- uma coesão ancestral que atravessa inúmeras comunidades, cada qual confluindo com a terra, com o tempo e entre si. Reler Antônio Bispo ensinando-me os modos de vida quilombolas faz ressoar em mim a mesma ancestralidade confluente entre os Innuit.\nota{Mauss e Beuchat, \enquote{Ensaio sobre as Variações Sazonais}; Santos, \textit{A Terra Dá, a Terra Quer}.}
% NOTA SUGERIDA [terminologia e comparação]: Conferir a glosa de Inuit e retirar a etimologia de “esquimó” enquanto não estiver documentada; o termo histórico pode permanecer no título da fonte. Delimitar os grupos e as fontes de Mauss e Beuchat e o contexto quilombola de Bispo, em vez de universalizar cada descrição.

Vamos pensar, por exemplo, na sensibilidade geracional de Innuits e quilombolas. Os primeiros mantêm  em suas coletividades  um mesmo conjunto de nomes pessoais, sendo as  crianças batizadas e incarnadas com o nome dos últimos ancestrais  falecidos. A geração neta revive a geração avó; para Bispo e os quilombolas, começo-meio-começo: geração avó, geração mãe, geração neta.
% NOTA SUGERIDA [localização]: Localizar nas duas obras as passagens sobre nomes de ancestrais e começo-meio-começo. Na passagem sobre sazonalidade, conferir “clãs” e quais agrupamentos o ensaio de fato descreve.

No inverno, estes se refugiam nas sombras das casas de farinha -- é o período do ano em que a vida coletiva é mais efervescente. Todos se encontram nas casas, trocam o que têm pelo que necessitam, ajudam com o trabalho caso não tenham nada, festejam, conversam; acima de tudo, trocam saberes, afetos e vivências, se envolvem entre si. No restante do ano, há sombras no mato e cada um vai trabalhar em sua própria roça  e a vida comunitária fica mais restrita ao núcleo familiar e à vizinhança.

Bispo relata uma vida social que respira junto com a terra nos quilombos; movimento que se encontra também na tradução por Mauss dos modos de vida Innuit; há em sua descrição duas configurações sociais que são igualmente determinadas pelas estações do ano: no verão, os diferentes núcleos familiares vivem isolados cada um em um abrigo próprio construído no início da estação, e pouco interagem com a vizinhança no cotidiano; contudo, no inverno, as famílias se unem em clãs e vivem todas num mesmo abrigo por clã, período em que as relações com a terra são mais amenas, e as do clã são intensas, havendo muitas festas, brincadeiras e compartilhamento de recursos durante o período.\nota{Mauss e Beuchat, \enquote{Ensaio sobre as Variações Sazonais}; Santos, \textit{A Terra Dá, a Terra Quer}.} O francês relata entre os Innuit uma respiração ainda mais acentuada do que aquela narrada por Bispo; ambas, no entanto, em profunda confluência com a terra.

Confluência que se avista também nas construções de cada povo. Ambos conversam intimamente com seu território, tecendo relações contratuais de profundo respeito e envolvimento. As construções Innuit variam, assim como as quilombolas, de acordo com sua localização; utilizam o que o território local tem a oferecer para construir casas e ferramentas atentas às condições de existência impostas e disponíveis localmente.

Penso-com as palavras de Bispo e Mauss em um saber ancestral comum a Innuits e quilombolas: o saber de confluir com a terra e todos os seus compartilhantes. E a ancestralidade me leva então para as últimas palavras de Ailton Krenak em \enquote{Futuro Ancestral} -- livro que rima na poesia, no espírito e na carne com o de Bispo, e termina dizendo: \enquote{o que nossas crianças aprendem desde cedo é a colocar o coração no ritmo da terra.}\nota{Krenak, \textit{Futuro Ancestral}, \enquote{O Coração no Ritmo da Terra}.} Ouso pensar-com Krenak que este saber ancestral indígena, que conflui com a ancestralidade de seus compartilhantes do extremo norte continental -- os Innuit -- e com a de seus compartilhantes quilombolas, constitui um saber ancestral de todos os povos da terra, de todos os seus compartilhantes, humanos ou não, e que se perdeu diante do monoteísmo cosmofóbico dos primeiros.
% NOTA SUGERIDA [alcance]: CONFERIR a página da citação direta na sua edição. Distinguir a ressonância construída pelo diário de uma origem ou de um saber universal demonstrados.

Retomando -- enfim -- o tema deste ensaio-estudo-diário não-moderno sobre como pensar o \enquote{Antropoceno}, penso que a Constituição moderna tecida por Latour rima com minha proposição dos antropocênicos e sua patologia comunicacional; mas rima ainda com um terceiro diagnóstico, o da cosmofobia enquanto patologia da ancestralidade de um povo que desaprendeu como respirar e colocar o coração no ritmo da terra, que se esqueceu deste saber e o pratica muito precariamente na vida higienizada e humanizada, desconectada da terra e de seus compartilhantes, que se enraizou nas cidades.

\entrada{19 de outubro de 2026}{Semeaduras}

Nego Bispo levou meu pensamento a semeaduras e com elas em mente comecei uma releitura do texto supracitado de Manuela Carneiro da Cunha. Antes de tratar da tripla rima com que encerrei minha entrada anterior, preciso tratar daquela entre Bispo e Manuela. Quando o leio dizendo que a agricultura quilombola não é mercadoria, remeto ao relato da antropóloga sobre a fala de um velho chefe yawanawa de que \enquote{Honi não é cultura!}\nota{Carneiro da Cunha, \enquote{\enquote{Cultura} e Cultura}.}
% NOTA SUGERIDA [citação direta]: CONFERIR na edição usada a formulação e a página do enunciado sobre honi. Nomear os interlocutores e o contexto conforme a fonte.

As articulações entre a crônica transmitológica de Nego Bispo e a crônica etnográfica de Manuela Carneiro já se iluminavam desde as primeiras páginas de \enquote{A terra dá, a terra quer}. \enquote{Precisamos transformar as armas do inimigo em nossa defesa}, diz o quilombola para introduzir a semeadura de palavras e sentidos na língua dos humanos como proteção contra o colonialismo destes.\nota{Santos, \textit{A Terra Dá, a Terra Quer}.} Bispo tem aí uma teccnologia ancestral de defesa interétnica ao avesso daquela descrita por Carneiro.
% NOTA SUGERIDA [citação direta]: CONFERIR a transcrição e a página da passagem de Bispo sobre transformar as armas do inimigo.

A antropóloga descreve o uso do estrangeirismo \enquote{cultura} por povos indígenas amazônicos, analisando mais especificamente a ocasião de uma negociação de direitos intelectuais sobre o uso em humanos do kampô --  uma toxina encontrada na pele de certas rãs. Embora possuam em suas línguas palavras para nomear seus costumes e modos de vida, a palavra \enquote{cultura} foi incorporada ao vocabulário de tais línguas; quando o velho chefe diz que honi não é \enquote{cultura}, ele não nega que o chá faça parte dos costumes e práticas que organizam e orientam a vida de seu povo. O estrangeirismo nomeia a \enquote{cultura} que os humanos colonizadores vão negociar, bloqueando os efeitos de reflexividade semântica que a tradução para termos nativos trariam.

Lembro de Latour dizendo que os pré-modernos misturavam as coisas naturais com as sociais. Isto é bem visível aqui, pois honi não é \enquote{cultura} -- mesmo sendo claramente um componente da cultura daquele chefe -- porque não antende às relações sociais articuladas em \enquote{cultura}. Objetos e sujeitos, natureza e sociedade, estão misturados nas palavras, de um modo que colabora para a manutenção das relações sociais sem enrijecê-las.

Na ocasião narrada pela antropólogca, o Kampô é \enquote{cultura}, mas honi não é. Isto lhe causa intriga. A diferenciação é imunização contra o colonialismo: porque nos imaginários humanos, os direitos intelectuais sobre objetos de cultura só podem ser individuais, coletivos ou públicos; nas relações com os povos da terra, isso restringe as negociações às formatações presentes nos imaginários metropolitanos; o uso de estrangeirismos em tais contextos interétnicos, além de melhor demarcar semanticamente o termo, protege as relações sociais incarnadas na língua nativa das contingências dos imaginários estrangeiros.\nota{Carneiro da Cunha, \enquote{\enquote{Cultura} e Cultura}. A leitura desse uso como imunização é uma hipótese deste diário.}

Honi não é \enquote{cultura} porque as relações sociais ancoradas neste termo estão limitadas pelo imaginário ocidental, e se articulam ali a uma noção de propriedade coletiva. Organizações indígenas aprenderam a falar a língua do colonizador para se protegerem de sua exploração, e a \enquote{cultura} é uma de suas principais defesas. Honi não é \enquote{cultura} porque, apesar de constituir a cultura nativa, não se articula a uma noção de propriedade coletiva, mas a um regime de conhecimento estranho ao imaginário ocidental, protegido de sua dominação mediante a diferenciação vocabular.

Se Nego Bispo semeia palavras e sentidos no vocabulário humano para proteger os modos de vida quilombola, Carneiro da Cunha relata um uso invertido do mesmo mecanismo. O velho chefe imunizava honi contra as sementes colonizadoras metropolitanas. Parafraseando Donna Haraway, importa muito mesmo quais nomes nomeiam as coisas.\nota{Haraway, \textit{Staying with the Trouble}, p. 12.}

Como eu havia dito ontem, encontrei em \enquote{A terra dá, a terra quer} uma encruzilhada entre os modos como os povos da terra se relacionam com ela e como se relacionam com os humanos. Bispo articulou em mim o primeiro elemento com o relato de Mauss sobre os Innuit, e o segundo com o de Manuela sobre a \enquote{cultura} entre povos indígenas no Brasil. Em ambos encontro uma sensibilidade fina às circunstâncias locais -- ao que se repete, e ao que se diferencia, em cada situação -- que fico tentado a situar como o saber ancestral mais fundamental: sempre escutar e aprender a escutar as diferenças --- as ancestrais e as emergentes.

\entrada{20 de outubro de 2026}{Uma patologia da comunicação}

Tem um ensaio do Latour -- \enquote{How to talk about the body?} -- com o qual comecei a pensar a proposta de Cecinepá como uma reação ao antropoceno. Este ensaio é muito importante para minha pesquisa desde a primeira vez que o li. Cecinepá é minha contribuição para as proposições ali apresentadas de composição progressiva e articulada de um mundo comum --- um multiverso --- e de ciências mais inter-essantes e inter-essadas.\nota{Latour, \enquote{How to Talk About the Body?}.}

Quero trabalhar hoje a noção de que o antropoceno é uma policrise derivada de uma patologia generalizada da comunicação: a patologia de não exergar que entre o sujeito e o objeto há uma forma que articula a emergência polarizada de ambos. O ensaio de Latour trabalha muito esta relação -- em minhas palavras -- entre sujeito, forma e objeto. Cecinepá é o componente mitológico da imunização que venho tentado compor afim de nos protegermos de tal patologia.
% NOTA SUGERIDA [promessa de desenvolvimento]: Se a nota apresentar Cecinepá, ligar mar das contingências, sol e ilhas de sal a um procedimento de comunicação; manter no corpo ao menos uma primeira apresentação antes do encerramento.

Latour assume o corpo como uma interface em constante processo de aprendizagem -- portanto, de transformação. Através desse processo, o corpo se torna progressivamente sensível a mais diferenças encontradas no mundo. A aprendizagem cultiva verdadeiros órgãos de sentido no corpo aprendiz; num só ato, ela faz emergirem um novo meio sensorial e um novo mundo sensível; transforma-se sujeito e objeto de uma só vez, articulados entre si.

Para ilustrar o processo, o francês apresenta o caso de perfumistas aprendizes e seu treinamento com kits de odores para se tornarem narizes. O kit-de-odores é um artefato construído por um especialista ao ordenar um certo número de pares de fragrâncias desde as mais sutilmente diferenciadas até as mais contrastantes entre si. Após algumas semanas expondo-se continuamente à sequência do kit, os corpos aprendizes progressivamente constituem narizes sensíveis a diferenças olfativas que até então lhes passavam despercebidas, mudas; transforma-se, assim, o sujeito, com um novo órgão olfativo; transforma-se, também, o mundo-objeto que ele habita, recheado agora de novos odores.\nota{Latour, \enquote{How to Talk About the Body?}, p. 206--207.}

O kit-de-odores e o processo de treinamento constituem a forma que proporciona a emergência deste corpo-mundo --- sujeito-objeto. O treinamento e o kit ordenam e propõem uma forma aos contrastes olfativos apresentados; articulam suas diferenças à forma da linearidade, da sequência linear, já sensível a corpos ocidentais. Não só aquelas diferenças mais sutis vão ficando mais marcantes por se situraem num lugar específico da sequência, mas as memórias que cada odor porventura invocar em cada corpo vão se alinhando também; esta sequência de odores e de contrastes é uma forma composta artificialmente que \guillemetleft ganha corpo\guillemetright ao ressoar uma forma linear já articulada pelo corpo aprendiz junto a uma sequência de diferenças sensíveis imperceptíveis a ele ecoando e formatando referenciais subjetivos particulares, circunstanciais a cada corpo; o kit-de-odores, enquanto uma metaforma, proporciona a sensibilização de um corpo-mundo a novos elementos que passam a compor um corpo`-mundo` -- a dupla articulada sujeito-objeto se modifica quando exposta continuamente a uma forma-sujeito articulada a uma forma-objeto, até que emerjam aderências entre a forma-sujeito e uma segunda forma-objeto, interna ao corpo; com isso, a forma-sujeito passa a ser localizada por duas formas-objeto, uma externamente proposta e outra internamente constituída; esta triangulação faz com que a forma-sujeito passe a ser ouvida como forma-objeto, de um modo que ainda verifica a consistência entre a articulação proposta entre as duas formas do kit e o mundo original habitado por aquele corpo; a forma -- o treinamento com o kit-de-odores -- cultiva a emergência de um novo par corpo-mundo, ou sujeito-objeto -- ou, ainda melhor, sujeito-oriente.
% NOTA SUGERIDA [divisão entre fonte e elaboração]: Transferir a triangulação e a metaforma para uma nota de desenvolvimento. No corpo, deixar o percurso kit → contrastes → corpo que aprende → mundo que se diferencia; linearidade, memórias e formas são hipóteses suas, não elementos a atribuir indistintamente a Latour.

Abolir o objeto em prol do oriente me soa bem por pelo menos dois grandes motivos: primeiro, objeto remete a manipulação, oriente remete a orientação, localização, o máximo de confiança que podemos esperar dos seres, das diferenças no mundo; segundo, o oriente contrapõe o ocidente, coletivo articulado à patologia à qual busco a imunização. Um terceiro motivo demorou um pouquinho a mais para chegar: há honra em ser um oriente para o outro, não há em ser um objeto.
% NOTA SUGERIDA [definição]: Experimentar “oriente” com um exemplo em que algo orienta sem ficar reduzido a recurso. Explicitar se a relação muda conforme quem se orienta e se há aqui um jogo vocabular ou uma oposição geográfica.

Rearticulando então: o kit-de-odores articula uma forma-sujeito e uma forma-oriente; através da exposição repetida a esta articulação, que chamarei de artifício, um corpo aprendiz progressivamente articula a forma-sujeito a uma segunda forma-oriente; aquela então se estabiliza enquanto forma-oriente; o corpo-mundo se atualiza em corpo`-mundo`. O artifício \guillemetleft ganhou corpo\guillemetright. Através da exposição continuada a um corpo, o artifício -- uma metaforma composta por uma forma-sujeito e uma forma-oriente -- se faz emergir como forma-oriente num corpo-mundo transformado; da metaforma germina um corpo-mundo transformado, a dualidade sujeito-oriente atualizada.

O sujeito é desorientador, forma sem forma, inconsistente; quando não, imperceptível. A forma-sujeito só aparece carregada por uma forma-oriente que aponte: aqui está uma forma que você não pode prever, ou mesmo perceber, a não ser que eu a indique. A familiaridade adquirida com um sujeito pode articulá-lo a algum outro oriente e transformá-lo, enfim, numa forma-oriente.

A patologia antropocênica decorre de uma elisão da forma enquanto mediação criativa do sujeito e do mundo que o orienta. Formas são os híbridos, laboratórios, experimentos, escolas, universidades, os artifícios diversos que se proliferam e somem, como se o mundo informasse de si mesmo. Perdemos a capacidade de comunicar -- de construir o comum -- porque perdemos as aspas, a capacidade de navegar e se articular entre diferentes regimes de conhecimento, narrações do mundo, mitologias, teorias, disciplinas, ciências, religiões, e por aí vai.

Nego Bispo denuncia o monoteísmo dos povos da cidade, e seus críticos acham que ele fala de religião. Monoteísmo é a cosmofobia, o medo do outro, a incapacidade de se articular a diferentes referenciais cosmológicos; ou seja, a diferentes deuses. O monoteísmo é doença que impede a comunicação e seu mecanismo pode ser descrito pela "Constituição dos modernos" transcrita por Latour.\nota{Santos, \textit{A Terra Dá, a Terra Quer}; Latour, \textit{Jamais Fomos Modernos}, cap. 2. A articulação entre cosmofobia e Constituição moderna é proposta pelo diário.}

Natureza transcendental, mas que pode ser agenciada como objeto para a transformação da realidade; sociedade imanente, que se assume anterior ao próprio sujeito; um abismo separando natureza e sociedade; e um deus barrado, juiz nada-poderoso ao qual se recorre nos períodos de crises e conflitos morais. Esta é a forma da elipse patológica, que se encontra mais nitidamente no signo enquanto categoria da linguagem, da comunicação e da semiótica: um significante e um significado, um sujeito e um objeto, um ativo e um passivo, separados, hierarquizados, o sujeito-sociedade-significante, o objeto-natureza-significado, com um deus barrado se escondendo no meio. O signo vem servindo como um cabresto que dificulta a articulação dos antropocênicos entre diferentes referenciais cosmológicos -- ou regimes de saber, ou \enquote{culturas}, ou \enquote{verdades}, ou teorias, ou ciências, ... -- e daí se dificulta progressivamente a comunicação enquanto composição de uma coletividade de diferentes em torno de um referencial comum, mas não único.
% NOTA SUGERIDA [contraponto]: Separar o esquema de garantias de Latour da sua analogia com o signo. Definir o que “forma” acrescenta à mediação e explicar por que a crítica recai sobre um modelo específico de signo.

\entrada{21 de outubro de 2026}{Forma}

Ando indeciso quanto à forma deste texto. Acho que eu gostaria que ele fosse dois: o texto + o meta-texto. Notas de rodapé robustas, como se o texto é que comentasse sobre elas, e não o contrário.

Iniciei este diário tendo em vista organizar minhas ideias para a escrita de um outro texto, que seria submetido para esta publicação. Logo notei o como o próprio diário era o texto que eu queria escrever, no formato que mais me apetece escrever, e assim segui; sem, no entanto, abandonar a imagem mental de um outro texto do qual este seria um comentário, um meta-texto. Eis que aparece o rodapé: o espaço perfeito no qual colocar o meu texto.

Escrevo, às avessas, notas de rodapé alçadas ao papel central, enquanto o protagonista vai para a coxia da página engrandecer-lhe. Comecei uns dois dias atrás a escrever pensando nas notas de rodapé que terei de fazer. Além de um diário de bordo, o que redijo é também uma carta de intenções, as quais venham a ser continuamente realizadas nos rodapés.
% NOTA SUGERIDA [nota que realiza]: Fazer este rodapé demonstrar o método: retomar uma intenção de 17 de outubro, explicitar a fonte, executar uma leitura e deixar uma pergunta que permita outra iteração.

Imagino uma prática científica orientada pela e para a comunicação. Um método que parte do discorrer de uma carta de intenções bem localizada e bem endereçada. Depende dessa carta e do público ao qual se dirige a escolha de referenciais teóricos, métodos, dados e experimentos e demais mediadores do trabalho.

A escrita, e mais especificamente a escrita linear, possui limitações comunicativas sérias. A oralidade tem uma força ancestral de trazer a palavra incarnada no presente, enquanto a escrita deixa a palavra imóvel. Quero aqui propor uma forma de escrita malemolente e que se faça presente, para que a palavra escrita dance como a palavra cantada.

Não serve aqui um texto \enquote{objetivo}, com sentido único e bem definido. Isso é uma falácia do signo. Palavras se movem a cada uso, e com sua ondulação semântica, vão emergindo dos textos novos sentidos.

Os rodapés estão aí para atualizações, correções e presentificações em geral. No método da carta de intenções, há esta vantagem de se poder ressituá-la em diferentes circunstâncias. Uma carta de intenções que aceite diferentes experimentos e teorias é mais articulada do que um artigo que só aceita replicações idênticas do procedimento original.

A carta de intenções dá à ciência o poder espiralar da oralidade. No discurso oral, cada performance carrega o que se repete junto do que se diferencia localmente. Junto ao rodapé bem localizado, a carta de intenções faz da ciência uma arte declamatória, um rodapé que presentifica a performance.
% NOTA SUGERIDA [tempo espiralar]: Se mobilizar a noção de tempo espiralar de Leda Maria Martins, distinguir a elaboração dela do seu uso para o diário e acrescentar a edição e a página efetivamente consultadas.

O leitor fica sempre convidado a refazer ou alterar os rodapés conforme a ocasião; conforme sua própria especulação das intenções propostas. Em seu ensaio, Latour se indaga sobre os aromas: estariam nas moléculas químicas? Ou nas memórias que perfumam? Com práticas científicas orientadas para a comunicação, uma mesma carta de intenções sobre possíveis benefícios da aromaterapia, por exemplo, poderia ter seus rodapés escritos por cientistas das mais diversas áreas, ou mesmo pelos narizes profissionais da indústria perfumista.

Qualquer pessoa pode, em última instância, fazer ciência declarativa, um formato científico open source iterate-to-evaluate --- ITE. Avalia-se uma carta de intenções iterando um novo texto de rodapé; avalia-se um texto de rodapé também com outro texto de rodapé. Com uma boa diagramação, uma carta pode até ser publicada com mais de um conjunto de textos de rodapé.
% NOTA SUGERIDA [avaliação]: Dar um exemplo de duas notas alternativas para a mesma intenção e indicar por quais critérios cada uma pode ser discutida: rastreabilidade, adequação à intenção e efeitos que a leitura torna observáveis.

\entrada{22 de outubro de 2026}{Cultura com aspas}

Todos os dias, antes de vir escrever aqui, tenho feito alguma leitura para me aquecer a mente e começar a movimentar as palavras (os totens) que a compõem. Hoje, dei continuidade ao ensaio de Manuela Carneiro da Cunha chamado \textit{\enquote{Cultura} e cultura: conhecimentos tradicionais e direitos intelectuais}. Quero pensar-com Manuela o \enquote{Antropoceno} com as aspas.\nota{Carneiro da Cunha, \enquote{\enquote{Cultura} e Cultura}.}

Em seu ensaio, a antropóloga distingue a cultura da \enquote{cultura} enquadrando-as em diferentes níveis lógicos. A segunda aparece como uma performance da primeira, e ainda, um metacomentário a seu respeito. A \enquote{cultura} como unidade elementar de um contexto interétnico compost por diferentes \enquote{culturas}; e a cultura como um agregado orgânico de práticas, costumes e saberes que organizam as diferenças (e são por elas organizadas) internamente a um coletivo de pessoas.

A \enquote{cultura} é uma performance cultural daquele coletivo frente a outros coletivos, constituindo uma rede de relações, não entre pessoas, mas entre grupos, entre clãs. E naquele contexto inter-\enquote{cultural}, os indivíduos são as próprias \enquote{culturas}, que se distribuem homogeneamente entre as pessoas que as compõem. Assim, na Copa do Mundo, por exemplo, surge o \enquote{brasileiro}, e cada brasileiro real compartilha com os demais a experiêcia da exibição de nossa performance \enquote{cultural} perante outras \enquote{culturas}, e se aproxima subjetivamente desse \enquote{brasileiro} que vai se produzindo, se realizando, em suas performances; agora, a três dias de uma votação que, independentemente de seu resultado, será um marco na história do Brasil, o país se divide entre dois clãs, numa eleição presidencial acirradíssima que pode constituir ordenamentos sociais profundamente distintos conforme o vencedor do pleito; na Copa, diante do Outro de fora, somos mais iguais; nas eleições, diante do Outro de dentro, somos mais diferentes; esta é a diferença entre estar no nível da \enquote{cultura} ou da cultura.
% NOTA SUGERIDA [contraexemplo]: Testar a afirmação de homogeneidade: quem pode falar como “brasileiro” e quem contesta a performance? O contraste Copa/eleição ilustra mudanças de enquadramento, mas não equivale ao contexto interétnico da fonte.

Quando proponho utilizar as aspas em \enquote{Antropoceno}, quero deslocar o discurso para o nível lógico de um contexto interétnico. Latour faz esse deslocamento em "Jamais fomos modernos" ao invocar os modernos e os não-modernos (entre outros) para a narrativa que tece; já em \enquote{How to talk about the body?}, ele opera de outro modo, focalizando não o corpo em si, mas \guillemetleft falas sobre o corpo\guillemetright, ou seja, falando de um \enquote{corpo} com aspas.\nota{Latour, \textit{Jamais Fomos Modernos}; Latour, \enquote{How to Talk About the Body?}. O uso das aspas para ler Latour é uma operação deste diário.} Quero mostrar a simetria entre o que fazem Bruno, Manuela, Ailton e Antônio, para então pensar o antropoceno a partir dessa simetria.

Sem aspas e com iniciação minúscula, o antropoceno é o tempo-espaço no qual vivem os antropocênicos; com aspas e em maiúscula, o \enquote{Antropoceno} se refere a \guillemetleft falas sobre o antropoceno\guillemetright. Minha proposta é de um \enquote{Antropoceno} no singular-processual -- tal qual o \enquote{corpo} de Latour, os rios para Krenak e Bispo, e a cultura (sem aspas) para Carneiro da Cunha -- que articule os diferentes discursos sobre o antropoceno. Proponho ainda que o \enquote{Antropoceno} se articula à patologia da comunicação que esconde a forma e apresenta os orientes como objetos.
% NOTA SUGERIDA [operação das aspas]: Distinguir a experiência histórica, o nome em disputa e a operação de segundo nível com que você observa seus usos. Mostrar o que muda numa frase ou numa prática ao passar de antropoceno a “Antropoceno”.

Essa escolha terminológica denuncia, a despeito da significação ou datação que se articule a ela, o pressuposto que é a grande divisão entre sujeito e \enquote{objeto} que a Constituição moderna operacionaliza. Penso o \enquote{Antropoceno} como a epítome desta Constituição. Se vivemos um período de múltiplas crises planetárias, é porque vemos problemas crescerem e se avolumarem à nossa frente e não encontramos agência para atravessá-los; é porque enquanto antropocênicos estamos subtraindo uma infinidade de agenciamentos possíveis e capazes de nos projetar para um exo-antropoceno, e então para um pós-antropoceno.

Sair do antropoceno demanda saltar do regime da cultura para o da \enquote{cultura}. Com esse deslocamento lógico-semiótico, enxerga-se o \enquote{Antropoceno} e se abre um imenso leque de oportunidades para ficar com o problema do antropoceno. Embarca-se, assim, no exo-antropoceno, um navio interétnico no qual não-antropocênicos fazem a travessia para o pós-antropoceno, quando os antropocênicos estarão extintos.
% NOTA SUGERIDA [figura de passagem]: Explicitar se a extinção nomeia o abandono de um modo de relação ou o desaparecimento de pessoas. Mostrar se exo- e pós-antropoceno são posições, tempos ou hipóteses; distinguir a travessia de uma sucessão inevitável.

O exo-antropoceno é o tempo-espaço quando não-antropocênicos constituem laços, saberes e práticas -- uma cultura exo-antropocênica -- para se imunizarem contra o \enquote{Antropoceno}. Ao longo da viagem, por vezes encontram antropocênicos à deriva no mar, sem uma forma na qual se agarrarem, e os resgatam; mas sabem que muitos se afogarão antes de conseguir embarcar; sabem também que, seja pelo afogamento, seja por embarcar no exo-antropoceno, o destino dos antropocênicos é a extinção. Apenas então poder-se-á aterrar na terra firme do pós-antropoceno.

\entrada{23 de outubro de 2026}{O piquenique cósmico}

Hoje li algo muito bonito do ensaio de Manuela. Ela falava sobre o conhecimento e seus diversos regimes. Disse que entre os povos amazônicos, a natureza é idiossincrática, e a cultura -- enquanto modo de conhecê-la, é universal; até mesmo os animais não-humanos compartilhariam com os humanos essa aculturação da natureza, apenas com seus referentes específicos.\nota{Carneiro da Cunha, \enquote{\enquote{Cultura} e Cultura}.}
% NOTA SUGERIDA [alcance etnográfico]: Localizar o trecho e distinguir a formulação da autora, o diálogo com o perspectivismo e sua leitura de “natureza idiossincrática”. Evitar atribuí-la indistintamente a todos os povos amazônicos.

O trecho já estava grifado, e certamente me ajudou a construir a imagem de um piquenique cósmico a orientar os encontros entre viventes. Imagino primeiramente um bosque -- essa natureza idiossincrática -- no qual os diferentes seres passeiam acompanhados de seus modos particulares de reconhecer as coisas do bosque e selecionar o que coletar em suas bolsas; num segundo momento -- do piquenique em si -- um conjunto de seres se reúne numa clareira no centro do bosque, estende uma toalha no chão, sentam-se, e compartilham suas coletas; reunidos, cada qual toma da toalha aquilo que lhes apetece, tudo que reconhecem como orientes para seus passeios pelo bosque. Penso em Antônio nos enquadrando como compartilhantes da terra; e num piquenique entre seres que compartilham seus diferentes modos de aculturar a terra, suas diferentes sensibilidades ao se relacionar com ela e com os demais compartilhantes.
% NOTA SUGERIDA [teste do piquenique]: Construir uma situação em que partilhar uma colheita não exija concordar sobre o que ela é. Perguntar quem escolhe a clareira, quem pode recusar e se os frutos chegam a ser reconhecíveis de ambos os lados.

Assim como não se pode conhecer a cultura do Outro imediatamente, apenas sua \enquote{cultura} -- performance da primeira num contexto intercultural --, um compartilhante não acompanha o modo de coletar dos demais, tendo acesso apenas ao resultado da colheita que estes compartilham consigo. O encontro com o Outro é sempre intercultura, daí a inacessibilidade à sua cultura, que fala através de performances específicas a cada encontro, constituindo máscaras -- tipo especial de oriente, espécie de persona que constela um conjunto de orientes (símbolos, práticas, técnicas, costumes) estabilizadas e continuamente re-estabilizadas nos encontros entre diferentes compartilhantes -- que realizam duplamente a comunicação entre eles. Por um lado, é mediante as máscaras que a comunicação acontece; por outro, elas são a própria comunalidade constituída em si entre diferentes compartilhantes que se comunicam.

No \enquote{Antropoceno}, a capacidade de constituir máscaras foi se esgotando com o monoteísmo cosmofóbico. Em vez de piqueniques, os antropocênicos fazem guerrars tacando os frutos de suas colheitas uns nos outros, pois são incapazes de enxergar na colheita alheia senão um veneno ao seu próprio modo de coletar. E não é de todo equivocada a visão, pois os piqueniques cósmicos são potentes imunizantes contra a cosmofobia antropocênica.

Se começarem a sentar-se à toalha com o Outro, os seres antropocênicos logo estariam fazendo colheitas muito mais ricas, comunicando com os demais compartilhantes modos muito mais agradáveis e seguros de viver-com a terra. Rapidamente se sintonizariam de volta ao ritmo desta provedora de tudo. Mas é preciso que vejam que o que tá parecendo veneno é vacina, que alcanem o exo-antropoceno para que enxerguem o \enquote{Antropoceno}.

\entrada{24 de outubro de 2026}{O paradoxo da cosmofobia}

Como dizer que o que parece uma coisa é outra sem soar colonialista, sem cair em cosmofobia? O piquenique é um encontro para trocas cosmológicas entre diferentes, mas aqueles que só sabem jogar sua colheita nos outros e que se nega a todo e qualquer fruto alheio não têm lugar nesta mesa. Seria este um caso análogo ao do paradoxo da tolerância?

Este diz que, para o convívio entre diferentes, é preciso recusar a presença do diferente intolerante, aquele que não tolera diferentes. O piquenique cósmico é uma heterotopia especulativa fundada não  na tolerância, mas na amizade e no apreço pelo diferente. Não seria de todo equivocado chamar ainda de uma experiência de amor ao próximo.
% NOTA SUGERIDA [alusão conceitual]: Identificar a edição de Popper e localizar a nota do paradoxo da tolerância. Sua formulação não se resume a excluir todo intolerante. Se “heterotopia” for conceito de Foucault aqui, localizar também a edição e explicar a transposição para a cena.

No piquenique, quanto mais diferenças se agregam e compartilham, mais amplo se torna o mundo de todos os presentes; mais articulado e com mais agenciamentos possíveis. Penso que o que permite superar o paradoxo é o seguinte: senta-se à mesa com quem se senta à mesa, joga-se frutos contra quem joga frutos; trata-se de reciprocidade -- com os cosmofóbicos também se aprende, seus frutos lançados também se aproveitam: eles ensinam a imunização contra aquilo que quer destruir o seu modo de existir. Ao Outro que teme a toalha no chão não se nega um lugar no banquete; tampouco se deixa de tomar para si os frutos arremessados que se reconheçam interessantes, nem de arremessar-lhes de volta os frutos por si próprio colhidos: isto se faz mostrando que a transformação dos modos de viver é o segredo para manter estes mesmos modos, e para uma boa vida com a terra.

Bispo conta que os povos indígenas ensinaram aos povos afroconfluentes que a ancestralidade que funcionava em sua terra nativa da qual foram retirados à força precisava mudar para funcionar do lado de cá do Atlântico.\nota{Santos, \textit{A Terra Dá, a Terra Quer}.} Houvessem então tradições afroconfluentes avessas à transformação, estas se perderam pois sua rigidez as enfraquece. No jogo cósmico da vida, a cosmofobia pode até render bons frutos localmente, no contexto da cultura, por algum período de tempo; no longo prazo, contudo, a terra seleciona os compartilhantes capazes de se transformarem junto a ela e os demais compartilhantes: aqueles que não temem o piquenique cósmico.
% NOTA SUGERIDA [inferência histórica]: Separar o relato de Bispo da hipótese do diário sobre tradições que se perderam; a segunda não decorre automaticamente do primeiro.

\entrada{25 de outubro de 2026}{Os antropocênicos e suas filas}

No momento em que cheguei na escola onde voto, tocou o sinal dando o aviso de que os portões deviam ser abertos. Diante de um espaço amplo, as pessoas faziam uma fila que se desfez assim que a larga passagem se abriu. Uma mulher deparou-se com dois policiais e os saudou: \enquote{que bom que estão aqui! Hoje tá complicado!}; \enquote{quê isso!?}, respondeu cordial e sorridente um dos soldados, \enquote{tá tudo tranquilo!}.

Fui o primeiro de minha seção a votar; isso porque não fui contaminado pela patologia da ordenação linear que acometeu os enfileirados que encontrei ao chegar ali. Foi bastante curioso e um tanto engraçado ver tanta gente numa linha reta aguardando a abertura de um portal pelo qual passariam seis ou sete pessoas ao mesmo tempo tranquilamente. Como não souberam confluir com o território, ordenaram-se sinteticamente numa linha reta desnecessária, ignorantes de que \enquote{tá tranquilo} se espalhar organicamente pela passagem.
% NOTA SUGERIDA [descrição e hipótese]: O que a fila permitia a quem não sabia como o portão abriria? Testar uma leitura em que a fila também coordena cuidado ou acesso; explicar quando a ordem se torna um obstáculo.

O saber sintético é assim: em prol de uma ordem pela ordem, condena à profanidade um amplo espaço de movimento e agência; porque não sabe confluir senão com uma noção abstraída de \enquote{ordem} que faz com que se façam filas desnecessárias por onde passariam multidões. O \enquote{Antropoceno} é, também, isto: esta patologia da mente linear. As múltiplas crises, como que uma só crise de confluência: os espaços de agência encontram-se interditados, profanados pela doença da linearidade -- a cosmofobia está matando os antropocênicos.

Sou otimista porque penso que a crise cosmofóbica vem se agudizando profundamente. Os entregadores de aplicativo conhecem bem a confluência dos cruzamentos frente a morosidade da linha reta. Se os carros se comportam como um grande rebanho sendo direcionado e redirecionado num complexo de currais, as motos são como cães confluindo com a terra e seus compartilhantes, correndo entre os bichos, se espalhando-com o território.

Sensibilidades para a confluência vêm emergindo e abrindo espaços. Vejamos o que se deu na campanha eleitoral para o segundo turno destas eleições presidenciais de 2026: metade do Brasil esteve em grande parte mobilizada por uma urgência existencial para construir confluências de toda ordem. Fazer campanhas políticas demanda mentalidades politeístas a fim de agenciar-se-com o Outro.

A vitória de Flávio Bolsonaro no primeiro turno das eleições foi um gatilho poderoso para a imunização contra a cosmofobia na esquerda brasileira. Buscou-se ocupar todos os espaços de agenciamento: desde estatísticas e dados de votação do TSE, até a articulação de aproximações cosmopolíticas com os mais diferentes segmentos em diversos eixos. Como falar com o agro, como falar com as mulheres, com os nêutros, a mobilização para dar caronas e assistências em geral para que as pessoas possam ir votar; num curto período de tempo, a busca urgente por \enquote{virar votos} fez com que a arbitragem do jogo fosse a capacidade para constituir comunalidades junto aos mais diferentes Outros, numa abertura cosmopolítica como pouco se vê entre os antropocênicos.
% NOTA SUGERIDA [documentação]: Usar “primeiro lugar no primeiro turno”, se confirmado. Acrescentar o resultado oficial desse turno e distinguir observação de campanha, dado e inferência sobre as alianças.

Em termos latourianos, os corpos brasileiros de esquerda tornaram-se bem mais articulados -- pois sensíveis a mais diferenças no mundo -- do que eram antes da corrida eleitoral.\nota{Latour, \enquote{How to Talk About the Body?}. O diagnóstico sobre a campanha é uma hipótese do diário, não uma conclusão da fonte.} Muito se aprendeu a ouvir o Outro, a confluir com os territórios -- especialmente o digital --, a olhar para o local que medeia o global. E independentemente de qual seja o resultado final, toda essa sensibilização já serviu para oxigenar a imaginação política contrariamente à linha reta, no sentido das confluências.

Pensar o \enquote{Antropoceno} articuladamente a uma patologia da comunicação abre espaços de agenciamento mediante articulações e cruzamentos ontoepistêmicos, cosmopolíticos. Desde o exo-antropoceno é mais fácil evitar as filas desnecessárias dos antropocênicos e confluir com os territórios, a terra, e seus diferentes compartilhantes. Uma prática científica politeísta é capaz de constituir comunicações mais potentes com os territórios, com compartilhantes humanos e não-humanos; uma cosmopolítica ativa e atenta é capaz de constituir arranjamentos coletivos, e intercoletivos, mais confluentes; uma ontologia politeísta é mais ágil na ginga necessária para sobreviver à queda do céu.
% NOTA SUGERIDA [alusão]: Se “queda do céu” remete a Davi Kopenawa e Bruce Albert, indicar edição e passagem; se não, decidir se a expressão deve continuar evocando essa obra.

\entrada{26 de outubro de 2026}{Um ponto de virada?}

Ontem foi um dia com o potencial de ser daqueles que encaminham os rumos históricos a renovadas direções. Pensar no sucesso da campanha pela eleição de Lula após a derrota para o candidato do fascismo no 1\textsuperscript{o} turno é, também, pensar no \enquote{Antropoceno}.\nota{Tribunal Superior Eleitoral, \enquote{Presidente do TSE Anuncia Reeleição de Lula}. A fonte confirma a apuração do segundo turno; a leitura política é do diário.} Porque não foram lulistas, petistas, progressistas, comunistas e eleitores de esquerda de modo geral os únicos que colaboraram com seus votos e sua voz para a \enquote{vitória da democracia}: o esforço reuniu pessoas da direita, ex-eleitores e ex-apoiadores de Jair Bolsonaro, e até mesmo a jornalista e ex-deputada Joice Hasselmann --- que foi a líder de seu governo na Câmara dos Deputados em 2019.
% NOTA SUGERIDA [caso de campanha]: Documentar a adesão de Joice Hasselmann e conferir a função exercida em 2019 (liderança do governo ou do Congresso). Separar a declaração pública, seu contexto e o efeito eleitoral, que o episódio isolado não mede.

Diante da extensiva patologia da comunicação que medeia e significa o \enquote{Antropoceno}, as últimas três semanas serviram como um tratamento de choque para a aprendizagem, a construção e a reprodução de vocabulários e de modos de fala e de escuta mais potentes para a constituição de comunalidades --- ou seja, para a confluência. Os diferentes grupos e pessoas de algum modo territorializados à esquerda tendem a ser, ironicamente, os mais adoecidos pela cosmofobia. É verdade que o fascismo se nutre dos nacionalismos, racismos e intolerâncias mais agúdos; mas ele sabe, como a direita sabe, agregar diferenças em torno das comunalidades que deseja constituir como ninguém.
% NOTA SUGERIDA [prova e escopo]: Escolher uma ou duas mediações de campanha observáveis. A passagem de uma cena a uma tese sobre “a esquerda” e “a direita” exige evidência e qualificação; considerar a eleição como coda ou como teste explícito da hipótese, se permanecer no ensaio.

Não se trata meramente de saber como falar, ou da vantagem política e comunicativa de já ter os referenciais de nossos imaginários e ideários marcados pelo status quo -- da velha questão marxiana entre estrutura e super-estrutura, ou base material e ideologia. O poder do capital e do dinheiro vem de sua profunda capacidade de se articularem e mediarem praticamente todas as experiências da vida de qualquer pessoa; o poder da comunicação de direita está na tomada dessa potência que atravessa todos os territórios e subjetividades atravessados por ela, partir de uma perspectiva que estabelece sem grandes esforços que ter mais dinheiro melhora a vida de qualquer um, e assim pode constituir uma legião de diferentes -- conservadores, reacionários, liberais, do povo mais pobre ao povo mais rico, nos campos e nas cidades, dos mais diferentes gêneros e orientações sexuais, de (quase) todos os tipos, enfim, de territórios e identidades -- que concordam que a liberdade para ganhar dinheiro é a única coisa que pode lhes pôr nas mãos a potencialidade de, por seu próprio esforço, melhorarem suas vidas e de suas famílias.

Nessa controvérsia, a esquerda tem a desvantagem de propor e lutar por contrapontos, amortizadores, quando não pela destituição completa desta força fundamental para a vida digna em virtualmente todos os territórios ocidentalizados. Mesmo em seus próprios territórios, o dinheiro e o capital aparecem como forças fundamentais; mas a comunicação costuma ser falha pois não se escuta o território do Outro atentamente para compreender como as políticas, propostas e linguagens esquerdistas muitas vezes lhe retira do único jogo que conhecem e que -- bem ou mal -- aprenderam a jogar para viver dignamente.

A centralidade da noção de \enquote{disputa}, de \enquote{competição}, também é outro grande entrave. É preciso aprender a negar a disputa como o único tabuleiro disponível para a dignidade. Do contrário, os contrapesos ao poder econômico e social como as políticas públicas de distribuição de renda, ou para as defesas das populações mais vulneráveis social e economicamente, serão traduzidas localmente nos territórios como privilégios e vantagens indevidas a classes específicas de jogadores enquadrados como concorrentes num mesmo jogo.


% Copie este bloco para acrescentar entradas:
% \entrada{[data]}{[título]}
%
% Primeiro período. Segundo período. Terceiro período.\nota{Desenvolvimento.}

% Lista para a seção Referências do projeto; ver instrução no cabeçalho.
% CONFERIR nas edições usadas as páginas pontuais das citações diretas.
% CONFERIR o intervalo de páginas do ensaio de Mauss e Beuchat na Ubu (2017).
% As referências a Popper, Foucault, Leda Maria Martins e à Queda do Céu
% estão sugeridas nos comentários; não entram na lista até a edição ser definida.
\newcommand{\referenciasdodiario}{%
  \referencia{Carneiro da Cunha, Manuela. \enquote{\enquote{Cultura} e Cultura: Conhecimentos Tradicionais e Direitos Intelectuais}. \textit{Cultura com Aspas e Outros Ensaios}, Cosac Naify, 2009, p. 311--373.}%
  \referencia{Crutzen, Paul J., e Eugene F. Stoermer. \enquote{The Anthropocene}. \textit{Global Change Newsletter}, n. 41, maio 2000, p. 17--18. \url{https://www.mpic.de/3864697/the-anthropocene}.}%
  \referencia{Haraway, Donna J. \textit{Staying with the Trouble: Making Kin in the Chthulucene}. Duke University Press, 2016. \url{https://doi.org/10.2307/j.ctv11cw25q}.}%
  \referencia{International Union of Geological Sciences e International Commission on Stratigraphy. \enquote{Joint Statement by the IUGS and ICS on the Vote by the ICS Subcommission on Quaternary Stratigraphy}. \textit{International Commission on Stratigraphy}, mar. 2024. \url{https://stratigraphy.org/news/152}.}%
  \referencia{Krenak, Ailton. \textit{Futuro Ancestral}. Companhia das Letras, 2022. \url{https://www.companhiadasletras.com.br/9786559211548-futuro-ancestral-3156/p}.}%
  \referencia{Latour, Bruno. \enquote{How to Talk About the Body? The Normative Dimension of Science Studies}. \textit{Body \& Society}, v. 10, n. 2--3, 2004, p. 205--229. \url{https://doi.org/10.1177/1357034X04042943}.}%
  \referencia{Latour, Bruno. \textit{Jamais Fomos Modernos: Ensaio de Antropologia Simétrica}. Traduzido por Carlos Irineu da Costa, Editora 34, 1994.}%
  \referencia{Mauss, Marcel, e Henri Beuchat. \enquote{Ensaio sobre as Variações Sazonais das Sociedades Esquimós: Estudo de Morfologia Social}. \textit{Sociologia e Antropologia}, de Marcel Mauss, traduzido por Paulo Neves, Ubu Editora, 2017.}%
  \referencia{Santos, Antônio Bispo dos. \textit{A Terra Dá, a Terra Quer}. Piseagrama e Ubu Editora, 2023.}%
  \referencia{Subcommission on Quaternary Stratigraphy. \enquote{Working Group on the \enquote{Anthropocene}}. \textit{International Commission on Stratigraphy}. \url{https://quaternary.stratigraphy.org/working-groups/anthropocene}.}%
  \referencia{Tribunal Superior Eleitoral. \enquote{Presidente do TSE Anuncia Reeleição de Lula e Faz Balanço do 2º Turno das Eleições 2026}. \textit{Tribunal Superior Eleitoral}, 26 out. 2026. \url{https://www.tse.jus.br/comunicacao/noticias/2026/Outubro/presidente-do-tse-anuncia-reeleicao-de-lula-e-faz-balanco-do-2o-turno-das-eleicoes-2026}.}%
}
